\documentclass[11pt]{article}
% ---------------------------------------------------------------
% R3.07 SQL dans un langage de programmation -- preambule commun des sujets de TP
% Article clair, fond blanc, accents de couleur coherents avec les slides
% Charte derivee de R5.A.08 Qualite de developpement (meme palette IUT)
% Utilisation :
%   \documentclass[11pt]{article}
%   % ---------------------------------------------------------------
% R3.07 SQL dans un langage de programmation -- preambule commun des sujets de TP
% Article clair, fond blanc, accents de couleur coherents avec les slides
% Charte derivee de R5.A.08 Qualite de developpement (meme palette IUT)
% Utilisation :
%   \documentclass[11pt]{article}
%   % ---------------------------------------------------------------
% R3.07 SQL dans un langage de programmation -- preambule commun des sujets de TP
% Article clair, fond blanc, accents de couleur coherents avec les slides
% Charte derivee de R5.A.08 Qualite de developpement (meme palette IUT)
% Utilisation :
%   \documentclass[11pt]{article}
%   \input{../../latex/common/r307-tp-preamble.tex}
%   \begin{document}
%   \rtpheader{1}{Titre du TP}
% ---------------------------------------------------------------
\usepackage[T1]{fontenc}
\usepackage[utf8]{inputenc}
\usepackage[french]{babel}
\usepackage[scale=0.92]{tgheros}
\usepackage{tgpagella}
\usepackage[a4paper, margin=2.2cm, headheight=14pt]{geometry}
\usepackage{graphicx}
\usepackage{xcolor}
\usepackage{listings}
\usepackage[skins,listings,breakable]{tcolorbox}
\usepackage{enumitem}
\usepackage{booktabs}
\usepackage{array}
\usepackage{multirow}
\usepackage{tikz}
\usetikzlibrary{shapes, shapes.geometric, arrows, arrows.meta, positioning, calc}
\usepackage[colorlinks=true, linkcolor=qaaccentdark, urlcolor=qaaccentdark]{hyperref}
\usepackage{fancyhdr}
\usepackage{lastpage}
\usepackage{titlesec}
\usepackage{etoolbox}

% ---------------------------------------------------------------
% PALETTE -- version imprimable des couleurs des diapositives
% ---------------------------------------------------------------
\definecolor{qaaccentdark}{HTML}{7A1680}   % violet du logo IUT, fonce pour rester lisible sur blanc
\definecolor{qaaccent2dark}{HTML}{9B4B94}  % violet clair du logo IUT (accent secondaire)
\definecolor{qapanellight}{HTML}{F4F6F8}   % fond de bloc tres clair
\definecolor{qalinelight}{HTML}{D8DEE4}    % filets
\definecolor{qatextdark}{HTML}{1A1D23}     % texte principal
\definecolor{qamuteddark}{HTML}{5B6472}    % texte secondaire
\definecolor{qacodekwl}{HTML}{7A1680}      % code : mots-cles
\definecolor{qacodestrl}{HTML}{2E7D32}     % code : chaines
\definecolor{qacodecoml}{HTML}{6B7280}     % code : commentaires
\definecolor{qacodeannl}{HTML}{B26A00}     % code : annotations / javadoc

\color{qatextdark}

\setlength{\parskip}{4pt}
\setlength{\parindent}{0pt}

% ---------------------------------------------------------------
% TITRES DE SECTION -- texte blanc sur bande de couleur pleine largeur
% (gere aussi les variantes etoilees \section*)
% ---------------------------------------------------------------
\newcommand{\qasectionband}[2]{%
  \noindent\colorbox{#1}{%
    \parbox[c]{\dimexpr\linewidth-2\fboxsep\relax}{%
      \color{white}#2%
    }%
  }%
}
\newcommand{\secformat}[1]{\qasectionband{qaaccentdark}{\bfseries\Large\thesection\quad #1}}
\newcommand{\secformatnn}[1]{\qasectionband{qaaccentdark}{\bfseries\Large #1}}
\newcommand{\subsecformat}[1]{\qasectionband{qaaccent2dark}{\bfseries\large\thesubsection\quad #1}}
\newcommand{\subsecformatnn}[1]{\qasectionband{qaaccent2dark}{\bfseries\large #1}}

\titleformat{\section}{\normalfont}{}{0pt}{\secformat}
\titleformat{name=\section,numberless}{\normalfont}{}{0pt}{\secformatnn}
\titleformat{\subsection}{\normalfont}{}{0pt}{\subsecformat}
\titleformat{name=\subsection,numberless}{\normalfont}{}{0pt}{\subsecformatnn}
\titleformat{\subsubsection}{\normalfont\bfseries\color{qaaccentdark}}{\thesubsubsection}{0.6em}{}
\titlespacing*{\section}{0pt}{14pt}{8pt}
\titlespacing*{\subsection}{0pt}{10pt}{6pt}
\titlespacing*{\subsubsection}{0pt}{8pt}{4pt}

% ---------------------------------------------------------------
% LISTINGS -- theme clair (lstlisting reste utilisable tel quel)
% \lstcode : police + couleur de base, a reutiliser dans les surcharges
% ---------------------------------------------------------------
\newcommand{\lstcode}{\ttfamily\color{qatextdark}}

\lstdefinelanguage{TypeScript}{
  keywords={typeof, new, true, false, catch, function, return, null, switch, var,
    if, in, while, do, else, case, break, const, let, interface, class, extends,
    implements, enum, export, import, this, super, constructor, public, private,
    protected, static, async, await, throw, throws, try, finally, void},
  ndkeywords={boolean, any, number, string, int, long, double},
  sensitive=true,
  comment=[l]{//},
  morecomment=[s]{/*}{*/},
  morestring=[b]',
  morestring=[b]"
}

\lstdefinestyle{r307code}{
  basicstyle=\lstcode\footnotesize,
  keywordstyle=\color{qacodekwl}\bfseries,
  ndkeywordstyle=\color{qaaccent2dark}\bfseries,
  stringstyle=\color{qacodestrl},
  commentstyle=\color{qacodecoml}\itshape,
  morecomment=[s][\color{qacodeannl}]{/**}{*/},
  backgroundcolor=\color{qapanellight},
  rulecolor=\color{qalinelight},
  frame=single,
  framesep=4pt,
  xleftmargin=6pt,
  xrightmargin=6pt,
  numbers=left,
  numberstyle=\tiny\color{qamuteddark},
  numbersep=8pt,
  keepspaces=true,
  columns=fullflexible,
  showstringspaces=false,
  showspaces=false,
  showtabs=false,
  breaklines=true,
  breakatwhitespace=false,
  captionpos=b,
  tabsize=2,
  escapeinside={(*@}{@*)},
}
\lstset{style=r307code}

% ---------------------------------------------------------------
% BOITES : exercice, code, information
% ---------------------------------------------------------------
\tcbset{
  exobox/.style={colback=qapanellight, colframe=qaaccentdark, boxrule=0.6pt,
    arc=2pt, left=8pt, right=8pt, top=6pt, bottom=6pt,
    coltitle=white, colbacktitle=qaaccentdark, fonttitle=\bfseries,
    boxsep=2pt, breakable},
  infobox/.style={colback=white, colframe=qaaccent2dark, boxrule=0.6pt,
    leftrule=3pt, arc=0pt, left=8pt, right=8pt, top=5pt, bottom=5pt,
    coltitle=qaaccent2dark, colbacktitle=white, fonttitle=\bfseries,
    breakable},
  codebox/.style={colback=qatextdark, colframe=qalinelight, boxrule=0.4pt,
    arc=2pt, left=8pt, right=8pt, top=5pt, bottom=5pt, listing only,
    listing options={basicstyle=\ttfamily\small\color{white},
      breaklines=true, showstringspaces=false, columns=fullflexible,
      keepspaces=true, frame=none, numbers=none, backgroundcolor={}}, breakable}
}
\newtcolorbox[auto counter]{exobox}[2][]{exobox,
  title={Exercice~\thetcbcounter\ifstrempty{#2}{}{ --- #2}}, #1}
\newtcolorbox{infobox}[1][\`A retenir]{infobox, title={#1}}
\newtcblisting{codebox}{codebox}

% ---------------------------------------------------------------
% EN-TETE DE PAGE (fil d'Ariane + pagination)
% ---------------------------------------------------------------
\newcommand{\qaheadtitlevalue}{}
\newcommand{\qaheadtitle}[1]{\renewcommand{\qaheadtitlevalue}{#1}}

\pagestyle{fancy}
\fancyhf{}
\renewcommand{\headrulewidth}{0.4pt}
\fancyhead[L]{\small\color{qamuteddark}R3.07 --- \qaheadtitlevalue}
\fancyhead[R]{\small\color{qamuteddark}\thepage\,/\,\pageref{LastPage}}
\fancyfoot[C]{\footnotesize\color{qamuteddark}IUT de Vannes --- Universit\'e Bretagne Sud}
\fancypagestyle{plain}{\fancyhf{}\renewcommand{\headrulewidth}{0pt}%
  \fancyfoot[C]{\footnotesize\color{qamuteddark}IUT de Vannes --- Universit\'e Bretagne Sud}}

% ---------------------------------------------------------------
% BANDEAU DE TITRE (pas une page dediee, pour rester compact)
% \rtpheader{numero ou intitule court}{titre du TP}
% ---------------------------------------------------------------
\newcommand{\ryear}{2026/2027}
\newcommand{\rtpheader}[2]{%
  \thispagestyle{plain}%
  \begin{minipage}[t]{0.68\textwidth}
    {\small\color{qamuteddark}
      Universit\'e Bretagne Sud\\
      IUT de Vannes\\
      D\'epartement Informatique --- \ryear}
  \end{minipage}%
  \hfill
  \begin{minipage}[t]{0.28\textwidth}
    \raggedleft
    \includegraphics[width=2.6cm]{iut_logo.jpg}
  \end{minipage}

  \vspace{6mm}
  \begin{center}
    {\large\bfseries\color{qatextdark} BUT 2 --- R3.07. SQL dans un langage de programmation}\\[2mm]
    {\normalsize\color{qamuteddark} TD/TP #1}\\[1mm]
    {\Large\bfseries\color{qaaccentdark} #2}
  \end{center}
  \vspace{4mm}
  \qaheadtitle{#2}
}

%   \begin{document}
%   \rtpheader{1}{Titre du TP}
% ---------------------------------------------------------------
\usepackage[T1]{fontenc}
\usepackage[utf8]{inputenc}
\usepackage[french]{babel}
\usepackage[scale=0.92]{tgheros}
\usepackage{tgpagella}
\usepackage[a4paper, margin=2.2cm, headheight=14pt]{geometry}
\usepackage{graphicx}
\usepackage{xcolor}
\usepackage{listings}
\usepackage[skins,listings,breakable]{tcolorbox}
\usepackage{enumitem}
\usepackage{booktabs}
\usepackage{array}
\usepackage{multirow}
\usepackage{tikz}
\usetikzlibrary{shapes, shapes.geometric, arrows, arrows.meta, positioning, calc}
\usepackage[colorlinks=true, linkcolor=qaaccentdark, urlcolor=qaaccentdark]{hyperref}
\usepackage{fancyhdr}
\usepackage{lastpage}
\usepackage{titlesec}
\usepackage{etoolbox}

% ---------------------------------------------------------------
% PALETTE -- version imprimable des couleurs des diapositives
% ---------------------------------------------------------------
\definecolor{qaaccentdark}{HTML}{7A1680}   % violet du logo IUT, fonce pour rester lisible sur blanc
\definecolor{qaaccent2dark}{HTML}{9B4B94}  % violet clair du logo IUT (accent secondaire)
\definecolor{qapanellight}{HTML}{F4F6F8}   % fond de bloc tres clair
\definecolor{qalinelight}{HTML}{D8DEE4}    % filets
\definecolor{qatextdark}{HTML}{1A1D23}     % texte principal
\definecolor{qamuteddark}{HTML}{5B6472}    % texte secondaire
\definecolor{qacodekwl}{HTML}{7A1680}      % code : mots-cles
\definecolor{qacodestrl}{HTML}{2E7D32}     % code : chaines
\definecolor{qacodecoml}{HTML}{6B7280}     % code : commentaires
\definecolor{qacodeannl}{HTML}{B26A00}     % code : annotations / javadoc

\color{qatextdark}

\setlength{\parskip}{4pt}
\setlength{\parindent}{0pt}

% ---------------------------------------------------------------
% TITRES DE SECTION -- texte blanc sur bande de couleur pleine largeur
% (gere aussi les variantes etoilees \section*)
% ---------------------------------------------------------------
\newcommand{\qasectionband}[2]{%
  \noindent\colorbox{#1}{%
    \parbox[c]{\dimexpr\linewidth-2\fboxsep\relax}{%
      \color{white}#2%
    }%
  }%
}
\newcommand{\secformat}[1]{\qasectionband{qaaccentdark}{\bfseries\Large\thesection\quad #1}}
\newcommand{\secformatnn}[1]{\qasectionband{qaaccentdark}{\bfseries\Large #1}}
\newcommand{\subsecformat}[1]{\qasectionband{qaaccent2dark}{\bfseries\large\thesubsection\quad #1}}
\newcommand{\subsecformatnn}[1]{\qasectionband{qaaccent2dark}{\bfseries\large #1}}

\titleformat{\section}{\normalfont}{}{0pt}{\secformat}
\titleformat{name=\section,numberless}{\normalfont}{}{0pt}{\secformatnn}
\titleformat{\subsection}{\normalfont}{}{0pt}{\subsecformat}
\titleformat{name=\subsection,numberless}{\normalfont}{}{0pt}{\subsecformatnn}
\titleformat{\subsubsection}{\normalfont\bfseries\color{qaaccentdark}}{\thesubsubsection}{0.6em}{}
\titlespacing*{\section}{0pt}{14pt}{8pt}
\titlespacing*{\subsection}{0pt}{10pt}{6pt}
\titlespacing*{\subsubsection}{0pt}{8pt}{4pt}

% ---------------------------------------------------------------
% LISTINGS -- theme clair (lstlisting reste utilisable tel quel)
% \lstcode : police + couleur de base, a reutiliser dans les surcharges
% ---------------------------------------------------------------
\newcommand{\lstcode}{\ttfamily\color{qatextdark}}

\lstdefinelanguage{TypeScript}{
  keywords={typeof, new, true, false, catch, function, return, null, switch, var,
    if, in, while, do, else, case, break, const, let, interface, class, extends,
    implements, enum, export, import, this, super, constructor, public, private,
    protected, static, async, await, throw, throws, try, finally, void},
  ndkeywords={boolean, any, number, string, int, long, double},
  sensitive=true,
  comment=[l]{//},
  morecomment=[s]{/*}{*/},
  morestring=[b]',
  morestring=[b]"
}

\lstdefinestyle{r307code}{
  basicstyle=\lstcode\footnotesize,
  keywordstyle=\color{qacodekwl}\bfseries,
  ndkeywordstyle=\color{qaaccent2dark}\bfseries,
  stringstyle=\color{qacodestrl},
  commentstyle=\color{qacodecoml}\itshape,
  morecomment=[s][\color{qacodeannl}]{/**}{*/},
  backgroundcolor=\color{qapanellight},
  rulecolor=\color{qalinelight},
  frame=single,
  framesep=4pt,
  xleftmargin=6pt,
  xrightmargin=6pt,
  numbers=left,
  numberstyle=\tiny\color{qamuteddark},
  numbersep=8pt,
  keepspaces=true,
  columns=fullflexible,
  showstringspaces=false,
  showspaces=false,
  showtabs=false,
  breaklines=true,
  breakatwhitespace=false,
  captionpos=b,
  tabsize=2,
  escapeinside={(*@}{@*)},
}
\lstset{style=r307code}

% ---------------------------------------------------------------
% BOITES : exercice, code, information
% ---------------------------------------------------------------
\tcbset{
  exobox/.style={colback=qapanellight, colframe=qaaccentdark, boxrule=0.6pt,
    arc=2pt, left=8pt, right=8pt, top=6pt, bottom=6pt,
    coltitle=white, colbacktitle=qaaccentdark, fonttitle=\bfseries,
    boxsep=2pt, breakable},
  infobox/.style={colback=white, colframe=qaaccent2dark, boxrule=0.6pt,
    leftrule=3pt, arc=0pt, left=8pt, right=8pt, top=5pt, bottom=5pt,
    coltitle=qaaccent2dark, colbacktitle=white, fonttitle=\bfseries,
    breakable},
  codebox/.style={colback=qatextdark, colframe=qalinelight, boxrule=0.4pt,
    arc=2pt, left=8pt, right=8pt, top=5pt, bottom=5pt, listing only,
    listing options={basicstyle=\ttfamily\small\color{white},
      breaklines=true, showstringspaces=false, columns=fullflexible,
      keepspaces=true, frame=none, numbers=none, backgroundcolor={}}, breakable}
}
\newtcolorbox[auto counter]{exobox}[2][]{exobox,
  title={Exercice~\thetcbcounter\ifstrempty{#2}{}{ --- #2}}, #1}
\newtcolorbox{infobox}[1][\`A retenir]{infobox, title={#1}}
\newtcblisting{codebox}{codebox}

% ---------------------------------------------------------------
% EN-TETE DE PAGE (fil d'Ariane + pagination)
% ---------------------------------------------------------------
\newcommand{\qaheadtitlevalue}{}
\newcommand{\qaheadtitle}[1]{\renewcommand{\qaheadtitlevalue}{#1}}

\pagestyle{fancy}
\fancyhf{}
\renewcommand{\headrulewidth}{0.4pt}
\fancyhead[L]{\small\color{qamuteddark}R3.07 --- \qaheadtitlevalue}
\fancyhead[R]{\small\color{qamuteddark}\thepage\,/\,\pageref{LastPage}}
\fancyfoot[C]{\footnotesize\color{qamuteddark}IUT de Vannes --- Universit\'e Bretagne Sud}
\fancypagestyle{plain}{\fancyhf{}\renewcommand{\headrulewidth}{0pt}%
  \fancyfoot[C]{\footnotesize\color{qamuteddark}IUT de Vannes --- Universit\'e Bretagne Sud}}

% ---------------------------------------------------------------
% BANDEAU DE TITRE (pas une page dediee, pour rester compact)
% \rtpheader{numero ou intitule court}{titre du TP}
% ---------------------------------------------------------------
\newcommand{\ryear}{2026/2027}
\newcommand{\rtpheader}[2]{%
  \thispagestyle{plain}%
  \begin{minipage}[t]{0.68\textwidth}
    {\small\color{qamuteddark}
      Universit\'e Bretagne Sud\\
      IUT de Vannes\\
      D\'epartement Informatique --- \ryear}
  \end{minipage}%
  \hfill
  \begin{minipage}[t]{0.28\textwidth}
    \raggedleft
    \includegraphics[width=2.6cm]{iut_logo.jpg}
  \end{minipage}

  \vspace{6mm}
  \begin{center}
    {\large\bfseries\color{qatextdark} BUT 2 --- R3.07. SQL dans un langage de programmation}\\[2mm]
    {\normalsize\color{qamuteddark} TD/TP #1}\\[1mm]
    {\Large\bfseries\color{qaaccentdark} #2}
  \end{center}
  \vspace{4mm}
  \qaheadtitle{#2}
}

%   \begin{document}
%   \rtpheader{1}{Titre du TP}
% ---------------------------------------------------------------
\usepackage[T1]{fontenc}
\usepackage[utf8]{inputenc}
\usepackage[french]{babel}
\usepackage[scale=0.92]{tgheros}
\usepackage{tgpagella}
\usepackage[a4paper, margin=2.2cm, headheight=14pt]{geometry}
\usepackage{graphicx}
\usepackage{xcolor}
\usepackage{listings}
\usepackage[skins,listings,breakable]{tcolorbox}
\usepackage{enumitem}
\usepackage{booktabs}
\usepackage{array}
\usepackage{multirow}
\usepackage{tikz}
\usetikzlibrary{shapes, shapes.geometric, arrows, arrows.meta, positioning, calc}
\usepackage[colorlinks=true, linkcolor=qaaccentdark, urlcolor=qaaccentdark]{hyperref}
\usepackage{fancyhdr}
\usepackage{lastpage}
\usepackage{titlesec}
\usepackage{etoolbox}

% ---------------------------------------------------------------
% PALETTE -- version imprimable des couleurs des diapositives
% ---------------------------------------------------------------
\definecolor{qaaccentdark}{HTML}{7A1680}   % violet du logo IUT, fonce pour rester lisible sur blanc
\definecolor{qaaccent2dark}{HTML}{9B4B94}  % violet clair du logo IUT (accent secondaire)
\definecolor{qapanellight}{HTML}{F4F6F8}   % fond de bloc tres clair
\definecolor{qalinelight}{HTML}{D8DEE4}    % filets
\definecolor{qatextdark}{HTML}{1A1D23}     % texte principal
\definecolor{qamuteddark}{HTML}{5B6472}    % texte secondaire
\definecolor{qacodekwl}{HTML}{7A1680}      % code : mots-cles
\definecolor{qacodestrl}{HTML}{2E7D32}     % code : chaines
\definecolor{qacodecoml}{HTML}{6B7280}     % code : commentaires
\definecolor{qacodeannl}{HTML}{B26A00}     % code : annotations / javadoc

\color{qatextdark}

\setlength{\parskip}{4pt}
\setlength{\parindent}{0pt}

% ---------------------------------------------------------------
% TITRES DE SECTION -- texte blanc sur bande de couleur pleine largeur
% (gere aussi les variantes etoilees \section*)
% ---------------------------------------------------------------
\newcommand{\qasectionband}[2]{%
  \noindent\colorbox{#1}{%
    \parbox[c]{\dimexpr\linewidth-2\fboxsep\relax}{%
      \color{white}#2%
    }%
  }%
}
\newcommand{\secformat}[1]{\qasectionband{qaaccentdark}{\bfseries\Large\thesection\quad #1}}
\newcommand{\secformatnn}[1]{\qasectionband{qaaccentdark}{\bfseries\Large #1}}
\newcommand{\subsecformat}[1]{\qasectionband{qaaccent2dark}{\bfseries\large\thesubsection\quad #1}}
\newcommand{\subsecformatnn}[1]{\qasectionband{qaaccent2dark}{\bfseries\large #1}}

\titleformat{\section}{\normalfont}{}{0pt}{\secformat}
\titleformat{name=\section,numberless}{\normalfont}{}{0pt}{\secformatnn}
\titleformat{\subsection}{\normalfont}{}{0pt}{\subsecformat}
\titleformat{name=\subsection,numberless}{\normalfont}{}{0pt}{\subsecformatnn}
\titleformat{\subsubsection}{\normalfont\bfseries\color{qaaccentdark}}{\thesubsubsection}{0.6em}{}
\titlespacing*{\section}{0pt}{14pt}{8pt}
\titlespacing*{\subsection}{0pt}{10pt}{6pt}
\titlespacing*{\subsubsection}{0pt}{8pt}{4pt}

% ---------------------------------------------------------------
% LISTINGS -- theme clair (lstlisting reste utilisable tel quel)
% \lstcode : police + couleur de base, a reutiliser dans les surcharges
% ---------------------------------------------------------------
\newcommand{\lstcode}{\ttfamily\color{qatextdark}}

\lstdefinelanguage{TypeScript}{
  keywords={typeof, new, true, false, catch, function, return, null, switch, var,
    if, in, while, do, else, case, break, const, let, interface, class, extends,
    implements, enum, export, import, this, super, constructor, public, private,
    protected, static, async, await, throw, throws, try, finally, void},
  ndkeywords={boolean, any, number, string, int, long, double},
  sensitive=true,
  comment=[l]{//},
  morecomment=[s]{/*}{*/},
  morestring=[b]',
  morestring=[b]"
}

\lstdefinestyle{r307code}{
  basicstyle=\lstcode\footnotesize,
  keywordstyle=\color{qacodekwl}\bfseries,
  ndkeywordstyle=\color{qaaccent2dark}\bfseries,
  stringstyle=\color{qacodestrl},
  commentstyle=\color{qacodecoml}\itshape,
  morecomment=[s][\color{qacodeannl}]{/**}{*/},
  backgroundcolor=\color{qapanellight},
  rulecolor=\color{qalinelight},
  frame=single,
  framesep=4pt,
  xleftmargin=6pt,
  xrightmargin=6pt,
  numbers=left,
  numberstyle=\tiny\color{qamuteddark},
  numbersep=8pt,
  keepspaces=true,
  columns=fullflexible,
  showstringspaces=false,
  showspaces=false,
  showtabs=false,
  breaklines=true,
  breakatwhitespace=false,
  captionpos=b,
  tabsize=2,
  escapeinside={(*@}{@*)},
}
\lstset{style=r307code}

% ---------------------------------------------------------------
% BOITES : exercice, code, information
% ---------------------------------------------------------------
\tcbset{
  exobox/.style={colback=qapanellight, colframe=qaaccentdark, boxrule=0.6pt,
    arc=2pt, left=8pt, right=8pt, top=6pt, bottom=6pt,
    coltitle=white, colbacktitle=qaaccentdark, fonttitle=\bfseries,
    boxsep=2pt, breakable},
  infobox/.style={colback=white, colframe=qaaccent2dark, boxrule=0.6pt,
    leftrule=3pt, arc=0pt, left=8pt, right=8pt, top=5pt, bottom=5pt,
    coltitle=qaaccent2dark, colbacktitle=white, fonttitle=\bfseries,
    breakable},
  codebox/.style={colback=qatextdark, colframe=qalinelight, boxrule=0.4pt,
    arc=2pt, left=8pt, right=8pt, top=5pt, bottom=5pt, listing only,
    listing options={basicstyle=\ttfamily\small\color{white},
      breaklines=true, showstringspaces=false, columns=fullflexible,
      keepspaces=true, frame=none, numbers=none, backgroundcolor={}}, breakable}
}
\newtcolorbox[auto counter]{exobox}[2][]{exobox,
  title={Exercice~\thetcbcounter\ifstrempty{#2}{}{ --- #2}}, #1}
\newtcolorbox{infobox}[1][\`A retenir]{infobox, title={#1}}
\newtcblisting{codebox}{codebox}

% ---------------------------------------------------------------
% EN-TETE DE PAGE (fil d'Ariane + pagination)
% ---------------------------------------------------------------
\newcommand{\qaheadtitlevalue}{}
\newcommand{\qaheadtitle}[1]{\renewcommand{\qaheadtitlevalue}{#1}}

\pagestyle{fancy}
\fancyhf{}
\renewcommand{\headrulewidth}{0.4pt}
\fancyhead[L]{\small\color{qamuteddark}R3.07 --- \qaheadtitlevalue}
\fancyhead[R]{\small\color{qamuteddark}\thepage\,/\,\pageref{LastPage}}
\fancyfoot[C]{\footnotesize\color{qamuteddark}IUT de Vannes --- Universit\'e Bretagne Sud}
\fancypagestyle{plain}{\fancyhf{}\renewcommand{\headrulewidth}{0pt}%
  \fancyfoot[C]{\footnotesize\color{qamuteddark}IUT de Vannes --- Universit\'e Bretagne Sud}}

% ---------------------------------------------------------------
% BANDEAU DE TITRE (pas une page dediee, pour rester compact)
% \rtpheader{numero ou intitule court}{titre du TP}
% ---------------------------------------------------------------
\newcommand{\ryear}{2026/2027}
\newcommand{\rtpheader}[2]{%
  \thispagestyle{plain}%
  \begin{minipage}[t]{0.68\textwidth}
    {\small\color{qamuteddark}
      Universit\'e Bretagne Sud\\
      IUT de Vannes\\
      D\'epartement Informatique --- \ryear}
  \end{minipage}%
  \hfill
  \begin{minipage}[t]{0.28\textwidth}
    \raggedleft
    \includegraphics[width=2.6cm]{iut_logo.jpg}
  \end{minipage}

  \vspace{6mm}
  \begin{center}
    {\large\bfseries\color{qatextdark} BUT 2 --- R3.07. SQL dans un langage de programmation}\\[2mm]
    {\normalsize\color{qamuteddark} TD/TP #1}\\[1mm]
    {\Large\bfseries\color{qaaccentdark} #2}
  \end{center}
  \vspace{4mm}
  \qaheadtitle{#2}
}

\graphicspath{{../../latex/common/assets/}{./}}
\lstset{language=SQL}

\begin{document}
\rtpheader{3}{Index et optimisation de requ\^etes}

\section{Objectifs et organisation}

Ce TP met en pratique le chapitre 3. Il dure environ \textbf{une heure} (apr\`es 30 minutes de cours).
Vous allez :
\begin{itemize}
  \item lire des plans d'ex\'ecution avec \texttt{EXPLAIN ANALYZE} ;
  \item cr\'eer des index et \textbf{mesurer} leur effet, en lecture comme en \'ecriture ;
  \item comprendre pourquoi un index n'est parfois pas utilis\'e ;
  \item optimiser des insertions depuis Java avec les lots (\textit{batch}).
\end{itemize}

\begin{infobox}[Ce qu'on attend]
Pour chaque exercice, notez dans un fichier \texttt{OPTIMISATION.md} de votre d\'ep\^ot :
la requ\^ete, le type de parcours lu dans le plan (\textit{Seq Scan}, \textit{Index Scan}\ldots) et le temps
d'ex\'ecution \textbf{avant} et \textbf{apr\`es} votre modification. Les temps varient d'une machine \`a l'autre :
c'est l'\textbf{ordre de grandeur} qui compte.
\end{infobox}

\section{Mise en place}

Le script \texttt{code/generer\_donnees.sql} (fourni sur Moodle) cr\'ee un sch\'ema s\'epar\'e \texttt{perf} avec
100\,000 clients, 200\,000 comptes et 1\,000\,000 d'op\'erations. Il ne touche pas \`a vos tables des TP pr\'ec\'edents.
\textbf{Seules les cl\'es primaires sont index\'ees.}

\begin{codebox}
docker cp generer_donnees.sql <conteneur>:/tmp/
docker exec -it <conteneur> psql -U user -d banque_db -f /tmp/generer_donnees.sql
docker exec -it <conteneur> psql -U user -d banque_db
\end{codebox}

Dans \texttt{psql}, placez-vous dans le sch\'ema et affichez le temps des requ\^etes :
\begin{lstlisting}
SET search_path TO perf;
\timing on
\d Operation
\end{lstlisting}

\section{Lire un plan d'ex\'ecution}

\begin{exobox}{Premiers plans}
\begin{enumerate}[label=\alph*)]
  \item Ex\'ecutez \texttt{EXPLAIN ANALYZE} sur les deux requ\^etes suivantes :
\begin{lstlisting}
SELECT * FROM Operation WHERE numCompte = 4242;

SELECT c.numClient, sum(o.montant)
  FROM Compte c JOIN Operation o ON o.numCompte = c.numCompte
 WHERE c.numClient = 4242
 GROUP BY c.numClient;
\end{lstlisting}
  \item Pour chacune : quel type de parcours est utilis\'e sur \texttt{Operation} ? Combien de lignes sont lues
        puis rejet\'ees (\texttt{Rows Removed by Filter}) ? Combien sont finalement retourn\'ees ?
  \item Comparez \texttt{EXPLAIN} et \texttt{EXPLAIN ANALYZE} sur la premi\`ere requ\^ete. Quelle est la diff\'erence ?
\end{enumerate}
\end{exobox}

\section{Indexer les cl\'es \'etrang\`eres}

\begin{exobox}{Index sur les cl\'es \'etrang\`eres}
\begin{enumerate}[label=\alph*)]
  \item Les colonnes \texttt{Operation.numCompte} et \texttt{Compte.numClient} sont des cl\'es \'etrang\`eres.
        Sont-elles index\'ees ? (\texttt{\textbackslash di} ou \texttt{SELECT * FROM pg\_indexes WHERE schemaname = 'perf';})
  \item Mesurez la suppression de 50 comptes, \textbf{dans une transaction annul\'ee} pour ne rien perdre :
\begin{lstlisting}
BEGIN;
DELETE FROM Operation WHERE numCompte BETWEEN 1 AND 50;
EXPLAIN ANALYZE DELETE FROM Compte WHERE numCompte BETWEEN 1 AND 50;
ROLLBACK;
\end{lstlisting}
        Pourquoi est-ce si long alors qu'on ne supprime que 50 lignes ? (Indice : que doit v\'erifier la base
        pour chaque compte supprim\'e ?)
  \item Cr\'eez les deux index :
\begin{lstlisting}
CREATE INDEX idx_operation_compte ON Operation(numCompte);
CREATE INDEX idx_compte_client    ON Compte(numClient);
\end{lstlisting}
        puis refaites les mesures des questions pr\'ec\'edentes. Quels types de parcours apparaissent maintenant ?
        Quel est le facteur de gain ?
\end{enumerate}
\end{exobox}

\section{Quand l'index n'est pas utilis\'e}

\begin{exobox}{Fonctions, \texttt{LIKE} et index composites}
\begin{enumerate}[label=\alph*)]
  \item Cr\'eez \texttt{CREATE INDEX idx\_client\_nom ON Client(nomClient);} puis analysez :
\begin{lstlisting}
SELECT * FROM Client WHERE nomClient = 'Nom4242';
SELECT * FROM Client WHERE lower(nomClient) = 'nom4242';
SELECT * FROM Client WHERE nomClient LIKE 'Nom424%';
\end{lstlisting}
        Lesquelles utilisent l'index ? Proposez et cr\'eez un index pour chacune des deux autres
        (voir le chapitre 3), puis v\'erifiez.
  \item Cr\'eez \texttt{CREATE INDEX idx\_op\_compte\_date ON Operation(numCompte, dateOperation);}.
        Pour chacune des requ\^etes suivantes, pr\'edisez d'abord si l'index sera utilis\'e, puis v\'erifiez :
\begin{lstlisting}
SELECT * FROM Operation WHERE numCompte = 42 AND dateOperation >= '2024-01-01';
SELECT * FROM Operation WHERE dateOperation = '2022-06-01';
\end{lstlisting}
  \item Cr\'eez un index sur \texttt{typeOperation} et analysez \texttt{SELECT * FROM Operation WHERE typeOperation = 'DEPOT';}.
        Quelle proportion des lignes est retourn\'ee ? Le gain est-il comparable \`a celui de l'exercice 2 ? Pourquoi ?
\end{enumerate}
\end{exobox}

\begin{exobox}{Le co\^ut des index}
\begin{enumerate}[label=\alph*)]
  \item Cr\'eez une copie vide de la table, sans index : \texttt{CREATE TABLE op\_copie (LIKE Operation);}
  \item Mesurez (\texttt{\textbackslash timing}) : \texttt{INSERT INTO op\_copie SELECT * FROM Operation LIMIT 300000;}
  \item Videz la table (\texttt{TRUNCATE op\_copie;}), cr\'eez-y 4 index (sur \texttt{numCompte}, \texttt{typeOperation},
        \texttt{dateOperation} et \texttt{(numCompte, dateOperation)}) et refaites la m\^eme insertion.
  \item Concluez : faut-il indexer \og toutes les colonnes, au cas o\`u \fg{} ?
\end{enumerate}
\end{exobox}

\section{Optimiser c\^ot\'e Java : les lots}

Dans votre projet du TP2, ajoutez une classe \texttt{BenchInsertions} qui ins\`ere \textbf{10\,000 op\'erations}
dans \texttt{perf.Operation} (num\'eros \`a partir de 2\,000\,001, \texttt{numCompte} entre 1 et 200\,000) et affiche la
dur\'ee de chaque variante mesur\'ee avec \texttt{System.nanoTime()}. Supprimez les lignes ins\'er\'ees entre deux variantes
(\texttt{DELETE FROM perf.Operation WHERE numOperation > 2000000}).

\begin{exobox}{Une ligne \`a la fois ou par lots ?}
\begin{enumerate}[label=\alph*)]
  \item \textbf{Variante 1} : un \texttt{PreparedStatement}, \texttt{executeUpdate()} pour chaque ligne, en
        \texttt{autoCommit = true}.
  \item \textbf{Variante 2} : m\^eme chose, mais dans \textbf{une seule transaction} (\texttt{setAutoCommit(false)} puis \texttt{commit()}).
  \item \textbf{Variante 3} : une seule transaction et \texttt{addBatch()} / \texttt{executeBatch()} par paquets de 1000.
  \item \textbf{Variante 4} : variante 3 en ajoutant \texttt{?reWriteBatchedInserts=true} \`a l'URL JDBC.
  \item Pr\'esentez les quatre dur\'ees dans un tableau et expliquez les \'ecarts (allers-retours r\'eseau,
        \'ecritures disque au \texttt{COMMIT}).
\end{enumerate}
\end{exobox}

\begin{lstlisting}[language=Java]
String url = "jdbc:postgresql://localhost:5432/banque_db";   // variante 4 : + "?reWriteBatchedInserts=true"
String sql = "INSERT INTO perf.Operation VALUES (?, ?, ?, ?, ?)";
long debut = System.nanoTime();
// ... insertions ...
System.out.printf("Variante %d : %d ms%n", v, (System.nanoTime() - debut) / 1_000_000);
\end{lstlisting}

\section{Pour aller plus loin (facultatif) : pagination}

\begin{exobox}{\texttt{OFFSET} ou pagination par cl\'e}
Comparez les plans et les temps de :
\begin{lstlisting}
SELECT * FROM Operation ORDER BY numOperation LIMIT 20 OFFSET 900000;
SELECT * FROM Operation WHERE numOperation > 900000 ORDER BY numOperation LIMIT 20;
\end{lstlisting}
Ajoutez \`a \texttt{BanqueDAO} une m\'ethode \texttt{List<Operation> pageOperations(int apresNum, int taille)} qui utilise la seconde forme.
\end{exobox}

\section{Rendu}
\begin{itemize}
  \item Le fichier \texttt{OPTIMISATION.md} avec vos mesures et vos explications.
  \item Un script \texttt{index.sql} contenant les index que vous recommandez pour l'application bancaire, chacun justifi\'e par
        une requ\^ete de l'application (en commentaire).
  \item La classe \texttt{BenchInsertions}.
\end{itemize}

\end{document}
